\section{Global X-Ray Properties}
\label{sec:global}

MACS~J0416.1-2403 has been previously identified as a merging galaxy cluster at $z=0.396$ \citep{Mann2012}. The brightness of the NE subcluster core is significantly more peaked than that of the SW subcluster, which instead has a rather flat brightness distribution (Figure~\ref{fig:xray-mosaic}). 

We determined the average cluster temperature by extracting spectra from a circular region of radius $1.2$~Mpc \citep[approximately $R_{\rm 500}$,][]{Sayers2013} centered at $\alpha=4^{\rm h}\, 16^{\rm m}\, 08.8^{\rm s}$ and $\delta=-24^{\circ}\, 04\arcmin \, 14.0\arcsec$ (J2000). The spectra extracted from the five \chandra\ ObsIDs were fit in parallel, assuming the cluster emission is described by a single-temperature thermal component\footnote{We tried adding an additional thermal component to fit the cluster emission, but the parameters of the second component could not be constrained.}. For this global fit, the sky background parameters were fixed to the values listed in Table~\ref{tab:bkg-model}, and the redshift was fixed to $0.396$. We found $T=10.06_{-0.49}^{+0.50}$~keV and $L_{\rm X, \,\, 0.1-2.4\,\,keV} = (9.14\pm 0.10)\times 10^{44}$~erg~s$^{-1}$. The average cluster parameters are shown in Table~\ref{tab:global}.

\input{global-model-tab}
